% META: {"id": "1NSI-WEB-IHM-RE-C1-CORRIGE", "chapitre": "1NSI-WEB-IHM", "type_objet": "corrige", "exercice_ref": "1NSI-WEB-IHM-RE-C1", "capacites": ["P-IHM-01A"], "capacites_codes": ["C1"], "mode_creation": "ex_nihilo", "sources_inspiration": [], "status": "needs_review"}
\begin{corrige}{1NSI-WEB-IHM-RE-C1}
\begin{enumerate}
  \item Prénom : \lstinline{<input type="text">}, l'utilisateur écrit.
        Ville imposée : \lstinline{<select>}, il choisit dans une liste fermée.
        Langues : des \lstinline{<input type="checkbox">}, plusieurs réponses possibles.
        Accord unique : une seule \lstinline{<input type="checkbox">}, ou une paire de boutons
        radio si l'on veut forcer un choix explicite.
  \item Aucun de ces quatre-là : c'est un \lstinline{<button type="submit">} qui déclenche
        l'envoi. Cocher, saisir ou choisir ne provoque aucune requête.
  \item Les boutons radio d'un même groupe sont \textbf{exclusifs} : en cocher un décoche
        l'autre. Or le visiteur peut parler plusieurs langues ; il faut donc des cases à
        cocher, qui sont indépendantes les unes des autres.
\end{enumerate}
\end{corrige}
